\section{Methodology}
\label{sec:methodology}

We present \textbf{RabbitVideo}, which exploits temporal coherence in diffusion models to accelerate sparse attention computation and enable memory-efficient inference on consumer GPUs. The system comprises three synergistic components: \textbf{Reuse-Across-Blocks} (RAB) for reducing K-means clustering overhead, \textbf{Block-Importance Tiers} (BIT) for hierarchical sparse attention computation, and \textbf{Attention-Informed Offloading} (AIO) for intelligent GPU memory management. Figure~\ref{fig:overview} illustrates the overall architecture.

\begin{figure}[t]
    \centering
    \includegraphics[width=\linewidth]{Figures/architecture.png}
    \caption{RabbitVideo overview. RabbitVideo amortizes semantic sparse attention across diffusion timesteps via a CPU-resident cluster cache (RAB), allocates attention computation using block-importance tiers (BIT), and fits within memory budgets by streaming transformer blocks with a sliding-window offloading runtime.}
    \label{fig:overview}
\end{figure}

\subsection{Preliminaries}
\label{sec:preliminaries}

\paragraph{Video Diffusion Transformers.}
Video diffusion models generate videos by iteratively denoising a latent representation $\mathbf{z}_t$ over $T$ timesteps. At each timestep $t$, a transformer processes the latent through self-attention layers:
\begin{equation}
    \text{Attention}(\mathbf{Q}, \mathbf{K}, \mathbf{V}) = \text{softmax}\left(\frac{\mathbf{Q}\mathbf{K}^\top}{\sqrt{d}}\right)\mathbf{V}
\end{equation}
where $\mathbf{Q}, \mathbf{K}, \mathbf{V} \in \mathbb{R}^{N \times d}$ are query, key, and value matrices with sequence length $N$ and head dimension $d$. For video generation, $N$ can exceed 100,000 tokens, making full attention with $\mathcal{O}(N^2)$ complexity prohibitively expensive.

\paragraph{Semantic-Aware Permutation (SAP).}
SAP~\cite{xi2025sparse,yang2025sparse} addresses this by clustering tokens using K-means and permuting them to achieve block-diagonal sparsity:
\begin{equation}
    \mathbf{Q}' = \mathbf{P}_Q \mathbf{Q}, \quad \mathbf{K}' = \mathbf{P}_K \mathbf{K}, \quad \mathbf{V}' = \mathbf{P}_K \mathbf{V}
\end{equation}
where $\mathbf{P}_Q$ and $\mathbf{P}_K$ are permutation matrices derived from K-means clustering with $C_Q$ query clusters and $C_K$ key clusters. A dynamic attention map $\mathbf{M} \in \{0,1\}^{C_Q \times C_K}$ determines which cluster pairs attend to each other, achieving sparse attention with complexity $\mathcal{O}(N \cdot B)$ where $B$ is the effective block size.

\paragraph{The Clustering Bottleneck.}
While SAP reduces attention complexity, K-means clustering introduces significant overhead. For each attention layer at each timestep, clustering requires:
\begin{equation}
    \mathcal{T}_{\text{kmeans}} = \mathcal{O}(N \cdot C \cdot I \cdot d)
\end{equation}
where $C$ is the number of clusters and $I$ is the number of iterations. With 60 layers and 50 timesteps, this results in 3,000 clustering operations per video generation.

\subsection{Key Insight: Temporal Coherence}
\label{sec:temporal_coherence}

Our key observation is that latent representations in diffusion models exhibit strong \textit{temporal coherence}—they change smoothly across consecutive timesteps. Formally, for latent representations $\mathbf{z}_t$ and $\mathbf{z}_{t-1}$:
\begin{equation}
    \|\mathbf{z}_t - \mathbf{z}_{t-1}\|_2 \ll \|\mathbf{z}_t\|_2
\end{equation}

This coherence manifests in three exploitable properties:
\begin{enumerate}
    \item \textbf{Cluster Stability}: Token cluster assignments remain largely stable across timesteps, with only boundary tokens shifting between clusters (exploited by RAB).
    \item \textbf{Attention Pattern Stability}: The importance distribution across cluster pairs changes gradually, allowing for hierarchical approximation (exploited by BIT).
    \item \textbf{Compute Cost Predictability}: Layer compute costs at timestep $t$ are predictive of costs at timestep $t+1$, enabling intelligent prefetching decisions (exploited by AIO).
\end{enumerate}

We exploit all three properties to amortize computation across timesteps, achieving substantial speedup without quality degradation.

\subsection{RAB: Reuse-Across-Blocks Cluster Cache}
\label{sec:ctca}

RAB reduces K-means overhead by reusing cluster assignments across timesteps. Instead of performing full clustering at every step, we maintain a cluster cache and selectively update it based on quality metrics.

\paragraph{Cluster Cache.}
For each attention layer $\ell$, we maintain a cache $\mathcal{C}^\ell$ containing:
\begin{equation}
    \mathcal{C}^\ell = \{(\mathbf{c}_Q, \mathbf{c}_K, \mathbf{a}_Q, \mathbf{a}_K, t_{\text{last}})\}
\end{equation}
where $\mathbf{c}_Q, \mathbf{c}_K$ are cluster centroids, $\mathbf{a}_Q, \mathbf{a}_K$ are token-to-cluster assignments, and $t_{\text{last}}$ is the timestep of last full clustering.

\paragraph{Adaptive Reclustering.}
At each timestep, we decide whether to perform full reclustering or reuse cached assignments based on two criteria:

\textit{Quality Criterion}: We measure cluster quality using the silhouette-like score:
\begin{equation}
    Q(\mathbf{a}, \mathbf{X}) = 1 - \frac{\sum_{i} \|\mathbf{x}_i - \mathbf{c}_{a_i}\|_2}{\sum_{i} \|\mathbf{x}_i - \bar{\mathbf{x}}\|_2}
\end{equation}
where $\mathbf{x}_i$ is the $i$-th token, $\mathbf{c}_{a_i}$ is its assigned centroid, and $\bar{\mathbf{x}}$ is the global mean. If $Q < \tau_q$ (quality threshold), we trigger reclustering.

\textit{Interval Criterion}: To prevent drift accumulation, we enforce reclustering after $\Delta_{\max}$ timesteps regardless of quality.

\paragraph{Centroid Update.}
When reusing assignments, we efficiently update centroids without full clustering:
\begin{equation}
    \mathbf{c}_k^{(t)} = \frac{1}{|S_k|} \sum_{i \in S_k} \mathbf{x}_i^{(t)}
\end{equation}
where $S_k = \{i : a_i = k\}$ is the set of tokens assigned to cluster $k$. This reduces complexity from $\mathcal{O}(NCI)$ to $\mathcal{O}(NC)$.

The complete RAB algorithm is presented in Algorithm~\ref{alg:ctca}.

\begin{algorithm}[t]
\caption{RAB: Reuse-Across-Blocks}
\label{alg:ctca}
\begin{algorithmic}[1]
\REQUIRE Tokens $\mathbf{X}$, cache $\mathcal{C}$, quality threshold $\tau_q$, max interval $\Delta_{\max}$
\ENSURE Cluster assignments $\mathbf{a}$, centroids $\mathbf{c}$

\STATE $\Delta \gets t_{\text{current}} - t_{\text{last}}$
\IF{$\mathcal{C}$ is empty \OR $\Delta \geq \Delta_{\max}$}
    \STATE $\mathbf{a}, \mathbf{c} \gets \text{KMeans}(\mathbf{X})$ \COMMENT{Full clustering}
    \STATE Update $\mathcal{C}$ with $(\mathbf{a}, \mathbf{c}, t_{\text{current}})$
\ELSE
    \STATE $\mathbf{c}' \gets \text{UpdateCentroids}(\mathbf{X}, \mathcal{C}.\mathbf{a})$
    \STATE $Q \gets \text{ComputeQuality}(\mathbf{X}, \mathbf{c}', \mathcal{C}.\mathbf{a})$
    \IF{$Q < \tau_q$}
        \STATE $\mathbf{a}, \mathbf{c} \gets \text{KMeans}(\mathbf{X})$ \COMMENT{Quality-triggered}
        \STATE Update $\mathcal{C}$
    \ELSE
        \STATE $\mathbf{a} \gets \mathcal{C}.\mathbf{a}$, $\mathbf{c} \gets \mathbf{c}'$ \COMMENT{Reuse}
    \ENDIF
\ENDIF
\RETURN $\mathbf{a}, \mathbf{c}$
\end{algorithmic}
\end{algorithm}

\subsection{BIT: Block-Importance Tiers}
\label{sec:ctaa}

While RAB accelerates clustering, the dominant cost in sparse attention is still the attention computation itself. BIT introduces a hierarchical attention mechanism that further exploits stability in block importance across timesteps.

\paragraph{Hierarchical Attention Tiers.}
We partition cluster pairs into three tiers based on their attention importance (Figure~\ref{fig:overview}):

\begin{enumerate}
    \item \textbf{Full Attention} ($\mathcal{F}$): High-importance pairs receive standard token-level attention.
    \item \textbf{Centroid Attention} ($\mathcal{A}$): Medium-importance pairs use centroid-to-centroid attention as approximation.
    \item \textbf{Skip} ($\mathcal{S}$): Low-importance pairs are skipped entirely.
\end{enumerate}

\paragraph{Importance Estimation.}
We estimate the importance of cluster pair $(i, j)$ using centroid dot products:
\begin{equation}
    w_{ij} = \frac{\exp(\mathbf{c}_Q^{(i)} \cdot \mathbf{c}_K^{(j)} / \sqrt{d})}{\sum_{j'} \exp(\mathbf{c}_Q^{(i)} \cdot \mathbf{c}_K^{(j')} / \sqrt{d})}
\end{equation}

For each query cluster $i$, we sort key clusters by importance and apply top-$p$ selection:
\begin{align}
    \mathcal{F}_i &= \{j : \sum_{j' \leq j} w_{ij'} \leq p_{\text{full}}\} \\
    \mathcal{A}_i &= \{j : p_{\text{full}} < \sum_{j' \leq j} w_{ij'} \leq p_{\text{total}}\} \\
    \mathcal{S}_i &= \{j : \sum_{j' \leq j} w_{ij'} > p_{\text{total}}\}
\end{align}

\paragraph{Centroid Attention Computation.}
For cluster pairs in $\mathcal{A}$, we compute attention at the centroid level:
\begin{equation}
    \mathbf{O}_{\text{centroid}}^{(i)} = \sum_{j \in \mathcal{A}_i} \alpha_{ij} \cdot \bar{\mathbf{V}}^{(j)}
\end{equation}
where $\bar{\mathbf{V}}^{(j)} = \frac{1}{|S_j|}\sum_{k \in S_j} \mathbf{v}_k$ is the mean value vector for cluster $j$, and $\alpha_{ij}$ is the centroid-level attention weight. This reduces complexity from $\mathcal{O}(|S_i| \cdot |S_j|)$ to $\mathcal{O}(1)$ per cluster pair.

\paragraph{Output Aggregation.}
The final output combines contributions from all tiers:
\begin{equation}
    \mathbf{O}^{(i)} = \mathbf{O}_{\text{full}}^{(i)} + \mathbf{O}_{\text{centroid}}^{(i)}
\end{equation}

The complete BIT procedure is detailed in Algorithm~\ref{alg:ctaa}.

\begin{algorithm}[t]
\caption{BIT: Block-Importance Tiers}
\label{alg:ctaa}
\begin{algorithmic}[1]
\REQUIRE $\mathbf{Q}, \mathbf{K}, \mathbf{V}$, assignments $\mathbf{a}_Q, \mathbf{a}_K$, thresholds $p_{\text{full}}, p_{\text{total}}$
\ENSURE Output $\mathbf{O}$

\STATE Compute centroids: $\mathbf{c}_Q, \mathbf{c}_K, \bar{\mathbf{V}}$
\STATE Compute importance: $\mathbf{W} = \text{softmax}(\mathbf{c}_Q \mathbf{c}_K^\top / \sqrt{d})$
\STATE Identify tiers: $\mathcal{F}, \mathcal{A}, \mathcal{S} \gets \text{TopPSelection}(\mathbf{W}, p_{\text{full}}, p_{\text{total}})$

\STATE $\mathbf{O} \gets \mathbf{0}$
\FOR{each query cluster $i$}
    \STATE \COMMENT{Full attention for high-importance pairs}
    \FOR{$j \in \mathcal{F}_i$}
        \STATE $\mathbf{O}[S_i] \mathrel{+}= \text{Attention}(\mathbf{Q}[S_i], \mathbf{K}[S_j], \mathbf{V}[S_j])$
    \ENDFOR
    \STATE \COMMENT{Centroid attention for medium-importance pairs}
    \STATE $\mathbf{O}[S_i] \mathrel{+}= \sum_{j \in \mathcal{A}_i} W_{ij} \cdot \bar{\mathbf{V}}^{(j)}$
\ENDFOR
\RETURN $\mathbf{O}$
\end{algorithmic}
\end{algorithm}

\subsection{Complexity Analysis}
\label{sec:complexity}

Table~\ref{tab:complexity} summarizes the computational complexity of each component.

\begin{table}[h]
\centering
\caption{Computational complexity comparison}
\label{tab:complexity}
\begin{tabular}{lcc}
\toprule
\textbf{Operation} & \textbf{Baseline} & \textbf{RabbitVideo} \\
\midrule
K-means per layer & $\mathcal{O}(NCI)$ & $\mathcal{O}(NC/R)$ \\
Attention (full) & $\mathcal{O}(N^2)$ & $\mathcal{O}(p_{\text{full}} \cdot N \cdot B)$ \\
Attention (centroid) & --- & $\mathcal{O}(C_Q \cdot C_K)$ \\
Memory transfer & Sequential & Overlapped (AIO) \\
\bottomrule
\end{tabular}
\end{table}

where $R$ is the measured RAB reuse ratio, $B$ is the average block size, and $p_{\text{full}}$ is a tunable threshold controlling the speed--quality trade-off.

\paragraph{Speedup Intuition.}
The combined speedup from RAB and BIT depends on the reuse ratio and the tier distribution. A simple estimate is:
\begin{equation}
    S_{\text{total}} = \frac{1}{p_{\text{full}} + p_{\text{centroid}} \cdot \epsilon + p_{\text{skip}} \cdot 0}
\end{equation}
where $\epsilon \ll 1$ is the relative cost of centroid attention. In practice we tune $(p_{\text{full}}, p_{\text{total}})$ to trace a speed--quality Pareto frontier and report the resulting runtime/quality trade-offs in Section~\ref{sec:evaluation}.

\subsection{AIO: Attention-Informed Offloading}
\label{sec:offloading}

To enable inference on memory-constrained GPUs (24GB), we introduce Attention-Informed Offloading (AIO), which uses signals from RAB and BIT to make intelligent prefetching and eviction decisions---going beyond naive FIFO scheduling.

\paragraph{Compute Cost Prediction.}
We predict the compute cost for layer $\ell$ at timestep $t$ using a weighted combination of signals:
\begin{equation}
    \hat{C}^{(\ell, t)} = w_\rho \cdot \rho^{(\ell, t-1)} + w_s \cdot s_{\text{rab}}^{(\ell, t-1)} + w_\tau \cdot \tau^{(\ell, t-1)}
\end{equation}
where $\rho$ is attention density from BIT (fraction of full-attention pairs), $s_{\text{rab}}$ is the RAB decision signal (1.0 for full recluster, 0.1 for reuse), and $\tau$ is normalized historical timing. Default weights are $(w_\rho, w_s, w_\tau) = (0.4, 0.3, 0.3)$.

\paragraph{Priority-Based Prefetching.}
Instead of sequential prefetching, we prioritize layers by predicted compute time:
\begin{equation}
    \text{Priority}(\ell) = \frac{1}{\text{dist}(\ell, \ell_{\text{current}})} \cdot (1 + 2\hat{C}^{(\ell)})
\end{equation}
Expensive layers are prefetched earlier, overlapping transfer with compute of current layers. This ensures high-cost layers are ready when needed.

\paragraph{Priority-Based Eviction.}
When GPU memory is full, we evict layers based on eviction score (lower = evict first):
\begin{equation}
    \text{EvictScore}(\ell) = 0.4 \cdot \text{Recency}(\ell) + 0.3 \cdot \hat{C}^{(\ell)} + 0.3 \cdot \text{Criticality}(\ell)
\end{equation}
This keeps expensive and soon-to-be-used layers on GPU, unlike FIFO which ignores compute cost.

\paragraph{Async Transfer with Pinned Memory.}
We use CUDA streams for asynchronous CPU-to-GPU transfer with pinned (page-locked) memory, achieving $\sim$2$\times$ faster transfers than pageable memory.

\begin{algorithm}[t]
\caption{AIO: Attention-Informed Offloading}
\label{alg:offload}
\begin{algorithmic}[1]
\REQUIRE Layer index $\ell$, GPU window size $W$, predictor $\mathcal{P}$
\ENSURE Layer $\ell$ ready on GPU

\IF{Layer $\ell$ not on GPU}
    \IF{Prefetch in progress for $\ell$}
        \STATE WaitForPrefetch($\ell$) \COMMENT{Prefetch hit}
    \ELSE
        \STATE SyncLoadToGPU($\ell$) \COMMENT{Prefetch miss}
    \ENDIF
\ENDIF

\STATE \COMMENT{Evict if GPU window full}
\IF{$|\text{LayersOnGPU}| \geq W$}
    \STATE $\text{victim} \gets \arg\min_{\ell'} \text{EvictScore}(\ell')$
    \STATE OffloadToCPU($\text{victim}$)
\ENDIF

\STATE \COMMENT{Priority-based prefetching}
\STATE $\text{targets} \gets \text{TopK}(\text{Priority}, k=2)$
\FOR{$t \in \text{targets}$}
    \STATE AsyncPrefetch($t$) \COMMENT{On prefetch stream}
\ENDFOR
\end{algorithmic}
\end{algorithm}

\paragraph{Signal Bridge.}
A lightweight signal bridge connects RAB and BIT to AIO: after each layer completes, the bridge records $(s_{\text{rab}}, \rho, \tau)$ for the predictor. This enables AIO to anticipate compute patterns without additional profiling passes.

\paragraph{Memory Configuration.}
We keep a sliding window of $W$ transformer blocks on GPU (default $W{=}6$ for 24GB), with remaining blocks on CPU in pinned memory. RAB centroids and assignments reside on CPU to minimize VRAM usage.

\subsection{Implementation Details}
\label{sec:implementation}

We implement RabbitVideo as drop-in replacements for standard attention in the HunyuanVideo pipeline using PyTorch and Triton. Key implementation choices:

\begin{itemize}
    \item \textbf{RAB Parameters}: Quality threshold $\tau_q = 0.8$, recluster interval $[\Delta_{\min}, \Delta_{\max}] = [2, 10]$ timesteps, adaptive mode enabled.
    \item \textbf{BIT Parameters}: $p_{\text{full}} = 0.90$--$0.95$, $p_{\text{total}} = 0.98$--$0.99$, minimum KC ratio $0.05$--$0.10$ for quality-speed tradeoff.
    \item \textbf{AIO Parameters}: GPU window $W = 6$ layers (24GB), prefetch count 2--4, prediction weights $(w_\rho, w_s, w_\tau) = (0.4, 0.3, 0.3)$, pinned memory enabled.
    \item \textbf{Warmup}: First 7.5\% of timesteps and 2.5\% of layers use full precision attention to ensure generation quality.
    \item \textbf{Hardware}: Tested on NVIDIA RTX 4090 (24GB) and A100 (40GB/80GB). RabbitVideo enables 720p, 5-second video generation on consumer GPUs that would otherwise require 40GB+ VRAM.
\end{itemize}
